\section{The marking towers and their determinant normalization}
\label{sec:towers}

The calculation uses two kinds of marking on the height-two stratum of a
height-three deformation.  A connected marking identifies its formal group
with a Honda group.  An \emph{integral} étale marking chooses a generator of
the Tate module of its étale quotient.  The distinction matters: the first
has structure group $J=\mathcal O_{D_2}^{\times}$, whereas the automorphism
group of the corresponding vector bundle is $D_2^{\times}$.
Theorem~\ref{thm:two-towers} gives the quotient comparison;
Lemmas~\ref{lem:etale-marking} and~\ref{lem:lift-torsors} retain its
normalized integral marking and identify the lift algebras.
In this section write $G=\mathcal G$ for the algebraic $p$-divisible
group fixed in Section~\ref{sec:setup}.

Put
\[
 d=p^2,\qquad q=p^3,\qquad A_2=A/(a)=k[[t]].
\]
We use Honda coordinates in which $[p]_{\Gamma_i}(T)=T^{p^i}$.
For an affinoid perfectoid $k$-algebra $(C,C^+)$, let
\[
 \mathcal H_i(C,C^+)=C^{\circ\circ},
\]
with addition given by $\Gamma_i$.  Since $C$ is perfect and $[p]$ is the
$p^i$-power map, $\mathcal H_i$ is a sheaf of $\Qp$-vector spaces.
All complements of zero in this section mean nonzero on every geometric
fiber.  The arguments below concern the characteristic-$p$ perfectoid
site over $k$.

\subsection{Finite torsion algebras}

We first retain the entire $p$-divisible group $G$, including the part that
becomes étale after $t$ is inverted.  Its finite kernels are formed over
$A$ by Weierstrass preparation and then base changed.  Thus, on $A_2$,
the algebra of $G[p^l]$ is the prepared finite algebra of $[p^l](T)$;
it has rank $p^{3l}$.  Formal completion at the identity is not taken
after passage to $K$.

\begin{lemma}[Primitive torsion coordinates]
\label{lem:primitive-coordinates}
For $l\geq 1$, let $\mathcal B_l$ be the reduced closure, over $A_2$,
of the points of $G[p^l]_K$ whose images generate the étale quotient.
Write $s_l$ for its point coordinate and set
$s_j=[p^{l-j}](s_l)$ for $1\leq j\leq l$.  Then
\begin{equation}
 \mathcal B_l=k[[s_l]],\qquad
 \mathcal C_l:=k[[t]][s_l^{d^l}]=k[[s_l^{d^l}]].
 \label{eq:primitive-rings}
\end{equation}
The algebra $\mathcal C_l[1/t]$ is the finite étale algebra of generators
of $G_K^{\mathrm{et}}[p^l]$.  Over it, the algebra of \emph{all} lifts of
the universal generator is $\mathcal B_l[1/t]$, finite free of rank $d^l$.
The transition maps are finite and injective, and
\begin{equation}
 s_j\in k[[s_l^{d^{l-j}}]]\qquad(l\geq j).
 \label{eq:primitive-transitions}
\end{equation}
\end{lemma}

\begin{proof}
Define the closure as the image of the finite $A_2$-algebra of $p^l$-torsion
in the reduced generic algebra of the indicated primitive locus.  It is
finite and reduced, every component dominates $\Spec A_2$, and its
closed fiber is supported at $s_l=0$.  It is therefore a complete local
ring with residue field $k$, dimension one, and maximal ideal $(t,s_l)$.
The element $s_1$ is nonzero on every generic component.

The series $[p](T)$ factors through $T^d$, and its reduction at $t=0$
is $T^q$.  Its coefficient of $T^d$ is $t$.  Dividing $[p](s_1)=0$
by $s_1^d$ in the reduced generic algebra gives
\[
 t\,u(s_1^d,t)+s_1^{q-d}=0,
 \qquad u(0,t)=1.
\]
Here $u$ is a unit power series.  This identity holds in the closure by
generic density.  It also holds in the complete finite subalgebra
$\mathcal C_l$: the series for $[p^{l-1}]$ factors through
$T^{d^{l-1}}$, so $s_1^d$ is a series in $t$ and $s_l^{d^l}$ with
zero constant term in the latter variable.  Since $q-d\geq d$, the
identity puts $t$ in the ideal $(s_l^{d^l})$ of $\mathcal C_l$.
Its maximal ideal is consequently generated by $s_l^{d^l}$.
Its dimension is one because it is integral over $A_2$.
The surjection
$k[[W]]\twoheadrightarrow\mathcal C_l$, $W\mapsto s_l^{d^l}$,
is an isomorphism: a nonzero ideal in $k[[W]]$ would lower dimension.
The same argument, now with generator $s_l$, proves the first identity
in \eqref{eq:primitive-rings}.  These are identifications of the
inclusion $\mathcal C_l\subset\mathcal B_l$, not just abstract
regularity assertions.

Write $[p^l](T)=g_l(T^{d^l})$ and prepare $g_l$.
The linear coefficient of $g_l$ is a unit multiple of
$t^{1+d+\cdots+d^{l-1}}$.  The invariant-differential identity for this
Frobenius-divided homomorphism shows that its derivative, in its prepared
kernel algebra, is that coefficient times a unit.  After inverting $t$
the prepared algebra is therefore étale, of rank $p^l$.
Its kernel before dividing by Frobenius has connected rank $d^l$.
The primitive images are distinguished exactly by $s_l^{d^l}$, so
their generator algebra is $\mathcal C_l[1/t]$.

The scheme of lifts of this generator is finite locally free of rank
$d^l$.  Its coordinate algebra surjects onto
$\mathcal B_l[1/t]$ by the same point coordinate.  But
\eqref{eq:primitive-rings} makes the latter free of rank $d^l$ over
$\mathcal C_l[1/t]$, with basis $1,s_l,\ldots,s_l^{d^l-1}$.
The surjection is an isomorphism.  In particular, reduction in the
definition of the closure has not removed any infinitesimal part of the
lift scheme on the exact-height locus.

Every primitive étale generator at level $j$ lifts at level $l$ after
a geometric extension.  The generic transition is consequently
surjective on spectra and injective on the reduced coordinate algebras;
generic density gives injectivity on the closures.  The transition is
finite because $\mathcal B_l$ is finite over $A_2$.  Finally,
$[p^{l-j}](T)$ factors through $T^{d^{l-j}}$.  Substituting the expression
for $t$ in $k[[s_l^{d^l}]]$ proves
\eqref{eq:primitive-transitions}.
\end{proof}

The connected marking algebra $L$ has a different construction.
Height classification gives an isomorphism of the formal law of
$G_K$ with $\Gamma_2$ over an algebraic closure
\cite[Theorem~IV]{Lazard55}.  In fact its coefficients are separable
over $K$.  Write the isomorphism as $f(T)=\sum_{n\geq1}b_nT^n$
and compare $f([p](T))=f(T)^d$.  At degree $nd$ this gives
\[
 b_n^d-t^n b_n=Q_n(b_1,\ldots,b_{n-1}),
\]
where $Q_n$ has coefficients in $K$.  This polynomial in $b_n$
has derivative $-t^n\ne0$.  The first equation is
$b_1^{d-1}=t$ with $b_1\ne0$.  Induction puts every coefficient
in a finite separable extension, so the entire isomorphism is
defined over a separable closure.  Its inverse has the same
property.  These successive finite separable coefficient equations,
with the formal-law identities imposed, construct the isomorphism
torsor by finite étale jet algebras.  Its automorphisms are the
constant Honda group $J$.  Thus
\[
 L=\colim_{U\subset J}L^U
\]
is an ind-étale $J$-torsor over $K$, with $U$ ranging over open subgroups.
This also gives the exact-height marking description recalled in
Section~\ref{sec:setup}; see the coheight-one marking
algebra in \cite[Section~2.7]{BMR26}.  Finite products of fields are
allowed at each stage.  The action of $P$ transports the deformation and
its connected marking, and the action of $J$ changes that marking.
They commute, fix $k$, and preserve the valuation normalized by $v(t)=1$.
In particular every $L^U$ is $P$-stable.  All actions on finite
coefficient data are continuous for the valuation topology.

\subsection{The normalized coordinate limit}

The following explicit comparison will also ensure that the frame
identification is integral.  Let $G_t$ be a specialization of the formal
law on $A_2$ to an affinoid perfectoid $k$-algebra $(C,C^+)$, with
$t\in C^{\circ\circ}$.  Write $P_t(T)=[p]_{G_t}(T)$.

\begin{lemma}[Universal-cover coordinates]
\label{lem:cover-coordinates}
There is a natural additive isomorphism
\[
 \Phi_t:\mathcal H_3(C,C^+)\xrightarrow{\ \sim\ }
 \varprojlim_{[p]}G_t(C^{\circ\circ}).
\]
Its coordinates and inverse are
\begin{equation}
 \Phi_t(z)_j=\lim_{r\longrightarrow\infty}
      P_t^{\circ r}\bigl(z^{1/q^{j+r}}\bigr),
 \qquad
 \Phi_t^{-1}((b_j))=\lim_{j\longrightarrow\infty}b_j^{q^j}.
 \label{eq:cover-limit}
\end{equation}
The isomorphism is $\Qp$-linear and is compatible with coefficient
maps and changes of deformation frame.  On an algebraically closed
perfectoid field, the zeroth projection is surjective and its kernel
is the integral Tate module of $G_t^{\mathrm{et}}$.
\end{lemma}

\begin{proof}
Use the spectral norm and choose $\lambda<1$ bounding $|t|$.
The coefficients of $P_t(T)-T^q$ and of $G_t-\Gamma_3$ have norm at
most $\lambda$.  Integral power series are $1$-Lipschitz on the open
unit ball, and factorization through $T^d$ gives
\begin{equation}
 |P_t^{\circ r}(v)-P_t^{\circ r}(w)|
       \leq |v-w|^{d^r}.
 \label{eq:cover-contraction}
\end{equation}
Successive approximations in the first formula in
\eqref{eq:cover-limit} differ by at most $\lambda^{d^r}$.
They therefore converge; their limits are topologically nilpotent and
satisfy
\[
 |\Phi_t(z)_j-z^{1/q^j}|\leq\lambda,
 \qquad P_t(\Phi_t(z)_{j+1})=\Phi_t(z)_j.
\]
For a compatible sequence $(b_j)$, consecutive normalized coordinates
differ by at most $\lambda^{q^j}$.  Hence its displayed inverse exists
and satisfies
\[
 |z-b_j^{q^j}|\leq\lambda^{q^j}.
\]
After taking the relevant root this error is at most $\lambda$;
applying \eqref{eq:cover-contraction} makes it tend to zero.
This proves one inverse identity.  Raising the first displayed
estimate to $q^j$ proves the other.  No common upper bound strictly
less than one for all the individual $|b_j|$ is required.

The coefficients of $\Gamma_3$ are fixed by $q$-power Frobenius.
The discrepancy between the specialized addition law and the Honda
addition law is bounded by $\lambda$; after normalization, or after
$r$ applications of $P_t$, it tends to zero by the same estimates.
Thus the formulas respect addition.  A change of deformation frame
differs from its Honda reduction by topologically small coefficients.
Its Honda reduction commutes with $q$-power Frobenius, since its
coefficients lie in $\mathbb F_{p^3}$.  The identical estimate proves
frame equivariance.  The maps commute with integral multiplication
and with inverse multiplication by $p$, so are $\Qp$-linear.
Continuity, including continuous scalar families, follows from uniform
convergence on smaller balls.

Over an algebraically closed field, preparation applied to
$P_t(T)-w$ gives a monic polynomial whose roots are small whenever
$w$ is small.  Choosing roots successively proves surjectivity of the
zeroth projection.  Its kernel consists of compatible $p^j$-torsion
points.  Connected finite groups have only the identity as a point
over a perfect field, so this kernel is precisely the Tate module of
the étale quotient.
\end{proof}

Let $X$ be the completed perfection of the punctured formal $t$-disc,
equivalently $\Spa(K^\flat)$.  Let $\mathcal T_{\mathrm{et}}\to X$ be the tower
of integral étale Tate generators.  By
Lemma~\ref{lem:primitive-coordinates}, it is the perfected inverse
tower of the primitive-coordinate spaces on $t\ne0$.

\begin{proposition}[Comparison of the entire generator tower]
\label{prop:generator-ball}
There is a natural $P\times\Zp^\times$-equivariant isomorphism
\begin{equation}
 \mathcal T_{\mathrm{et}}\simeq\mathcal H_3\setminus\{0\},
 \qquad (s_l)\longmapsto z=\lim_l s_l^{q^l}.
 \label{eq:generator-ball}
\end{equation}
It identifies completed function algebras, with their spectral norms,
on an exhaustion by smaller closed balls.  The assertion commutes
with base change and remains valid after adjoining any compact
profinite parameter space.
\end{proposition}

\begin{proof}
Retain the reduced closures from
Lemma~\ref{lem:primitive-coordinates}, including their point at $t=0$.
After perfection, a point of the étale quotient has a unique lift
through the connected factor.  The resulting inverse tower is thus
the tower of the rings $\mathcal B_l$.

We check its complete algebras after extension to an algebraically
closed perfectoid field $C$ of characteristic $p$.  Fix
$0<\rho<1$ in its value group and put $y_l=s_l^{q^l}$.
At level $l$ take the closed disc $|y_l|\leq\rho$.
At a point of this disc, write $\lambda=|t|$.
The proof of Lemma~\ref{lem:primitive-coordinates} at level one,
and finite telescoping of normalized coordinates, give
\[
 \lambda\leq |s_1|^d,
 \qquad |s_1|\leq\max(\rho^{1/q},\lambda).
\]
If $\lambda>\rho^{1/q}$ these inequalities would imply
$0<\lambda\leq\lambda^d<\lambda$.  Consequently
\begin{equation}
 \lambda\leq\rho^{d/q},\qquad
 |y_{l+1}-y_l|\leq\lambda^{q^l}
                   \leq\rho^{d q^{l-1}}<\rho.
 \label{eq:generator-radius}
\end{equation}
The finite transition maps preserve these discs.  Conversely, every
geometric point of a level-$l$ disc has a lift by finite surjectivity,
and the strict inequality in \eqref{eq:generator-radius} forces its
lifts to remain in the next disc.  Pullback is therefore isometric
for the spectral norm, also after completed perfection.

In the completion of the direct limit of these perfected Gauss
algebras, $y_l$ converges to $z$.  Send the coordinate of the perfected
radius-$\rho$ disc to $z$.  The resulting homomorphism is isometric:
for a polynomial in finitely many fractional powers, evaluation at
$y_l$ has exactly its Gauss norm, since $s_l$ is a free disc
coordinate and $q^l$ is a power of $p$.  These evaluations converge
to evaluation at $z$.  For a nonzero polynomial, their difference
is eventually smaller than its fixed norm, so the limiting norm is
the same.  Completion preserves this isometry.

The image is dense.  For $l\geq j$, write
$s_j=F_{j,l}(s_l^{d^{l-j}})$ as in
\eqref{eq:primitive-transitions}, with $F_{j,l}\in k[[T]]$.
Replacing $s_l$ by $z^{1/q^l}$ changes this expression by at most
\[
 \lambda^{d^{l-j}}\leq\rho^{(d/q)d^{l-j}},
\]
which tends to zero with $l$.  The substituted series converges on
the $z$-disc.  The same approximation after taking roots proves
that every $s_j$ and each of its $p$-power roots belongs to the
closed image.  These generate the completed limit algebra.
The isometry is therefore onto.

Here is the relative descent explicitly.  Let
$E=\widehat{k((u))^{\mathrm{perf}}}$, with $u$ an auxiliary
variable, and choose $C=\widehat{\overline E}$.
For an affinoid perfectoid test $S=\Spa(A,A^+)$, a
pseudouniformizer $\pi\in A^+$ and its unique compatible roots
give a continuous map $E\to A$, $u\mapsto\pi$.
The field extension $\Spa(C,C^\circ)\to\Spa(E,E^\circ)$ is
a $v$-cover, so its base change
\[
 S_C=S\times_{\Spa(E,E^\circ)}\Spa(C,C^\circ)\longrightarrow S
\]
is a $v$-cover.  The calculation over $C$ established an
isomorphism of complete algebras, hence of perfectoid spaces,
on every smaller disc.  It therefore applies to all $C$-tests,
including $S_C$, rather than only to geometric points.
The coordinate map is defined over $k$; its inverse over this
cover is unique and agrees on the \v Cech overlap, so it descends
to $S$.  This uses base change of an established isomorphism,
not preservation of a spectral isometry by an arbitrary
completed tensor product.  Smaller discs exhaust the comparison:
on a quasicompact test $t$ is uniformly small, and the first
torsion coordinate and \eqref{eq:generator-radius} bound all
normalized coordinates uniformly away from radius one.

It remains to identify the puncture.  If $z=0$ on a geometric fiber,
Lemma~\ref{lem:cover-coordinates} gives $|s_l|\leq\lambda$ for all
$l$.  For fixed $j$,
\[
 |s_j|=|P_t^{\circ(l-j)}(s_l)|
           \leq\lambda^{d^{l-j}}\longrightarrow0.
\]
Thus all $s_j$ vanish, and the primitive-coordinate equation forces
$t=0$.  Conversely, at $t=0$ the reduced torsion coordinates vanish.
Removing these corresponding zero loci gives
\eqref{eq:generator-ball}.  Equivariance is the frame and scalar
equivariance of Lemma~\ref{lem:cover-coordinates}.
Every estimate used a uniform spectral norm; taking the supremum
over a compact profinite parameter preserves each estimate and each
density argument.
\end{proof}

\subsection{Relative bundles and integral frames}

We spell out the vector-bundle results used in the comparison.
For a perfectoid test $S$, write $X_S$ for its relative
Fargues--Fontaine curve.  Relative cohomology means the associated
sheaf on the $v$-site; it does not posit a morphism of schemes
$X_S\to S$.  The bundles $V_i$ are the Honda forms of rank $i$
and degree one.

\begin{lemma}[Relative section facts]
\label{lem:relative-section-facts}
There are natural isomorphisms of additive sheaves
\[
 \mathcal H_i\simeq \mathcal H^0(V_i),\qquad
 \mathcal H^0(\mathcal O)=\underline{\Qp}.
\]
The automorphism sheaf of $V_i$ is the locally profinite group
$D_i^\times$.  A family fiberwise isomorphic to $V_i$ is locally
of that form on the perfectoid site.  Trivial bundles and these
positive basic bundles have no higher relative cohomology sheaves.
For such bundles, $\underline{\Qp}$-linear maps between the section
sheaves are exactly the maps of bundles, relatively and compatibly
with base change.
\end{lemma}

\begin{proof}
The Lubin--Tate universal-cover description of sections of
$\mathcal O(1)$ is
\cite[Proposition~II.2.2]{FS21}.  Applying it with the unramified
coefficient field of degree $i$ and pushing forward gives
$\mathcal O(1/i)=V_i$ and the Honda height-$i$ cover.
The positive-section calculation, including the perfected-ball
description, is also
\cite[Proposition~II.2.5]{FS21}.
The classification on geometric curves and the relative pure-family
theorem give the stated local forms and automorphisms; see
\cite[Theorems~II.2.14 and II.2.19]{FS21}.
Equivalently, the basic stratum is the classifying stack of the
corresponding division-algebra group
\cite[Theorem~III.4.5]{FS21}.
The Honda forms are defined over our $k$: commuting with
$[p](T)=T^{p^i}$ puts every coefficient of a geometric Honda
endomorphism in $\mathbb F_{p^i}\subset k$.

For precision about maps, let $\tau$ denote the morphism of sites
used for relative curve cohomology.  Ansch\"utz--Le~Bras prove that
\[
 R\tau_*:\Perf(X_S)\longrightarrow
 D(S_v,\underline{\Qp})
\]
is fully faithful for every small $v$-stack $S$
\cite[Corollary~3.11]{ALB25}.
The same section calculation gives
$R\tau_*\mathcal O=\underline{\Qp}$ and
$R\tau_*V_i=\tau_*V_i[0]$; the latter uses positivity.
Taking degree-zero morphisms in this fully faithful functor proves
the last assertion.  These statements are relative on $S$, so they
apply over the finite-field base site and to forms obtained by
descent.  In particular, the assertion about $\mathcal O$ is
vanishing of its higher \emph{relative sheaves}, not an assertion
of vanishing on every unrefined test object.
\end{proof}

\begin{lemma}[The two stable-bundle tests]
\label{lem:stable-bundle-tests}
A section of $V_3$ nonzero on every geometric fiber gives a
subbundle $\mathcal O\hookrightarrow V_3$ with quotient locally
isomorphic to $V_2$.  Conversely, an extension
\begin{equation}
 0\longrightarrow\mathcal O\longrightarrow E
       \longrightarrow V_2\longrightarrow0
 \label{eq:basic-extension}
\end{equation}
has middle bundle fiberwise $V_3$ if and only if its extension class
is nonzero on every geometric fiber.  Moreover,
$\mathcal H^0(V_2^\vee)=0$.
\end{lemma}

\begin{proof}
Work first on a geometric curve.  The saturation of the image of a
nonzero section $\mathcal O\to V_3$ is a line bundle of nonnegative
integral degree.  Stability of $V_3$ bounds that degree by $1/3$,
so it is zero.  The section has no zeros, since a nonzero effective
divisor has positive degree.  Its quotient has rank two and degree
one.  A quotient of that quotient of nonpositive slope would also
be a quotient of $V_3$, contrary to semistability.  Its slopes are
therefore positive.  The classification of vector bundles, and
integrality of the degree of each stable summand, force the
quotient to be $V_2$.

An extension \eqref{eq:basic-extension} has no negative-slope
quotient, since neither end has a nonzero map to a negative-slope
bundle.  If it has a slope-zero quotient, that quotient receives
a nonzero map from the kernel $\mathcal O$, because
$\Hom(V_2,\mathcal O)=0$.  The quotient is a sum of copies of
$\mathcal O$; compose with a linear functional nonzero on that
kernel image and rescale.  This gives a retraction onto the
kernel, so the extension splits.  A nonsplit
extension thus has only positive slopes.  A positive bundle of
rank three and degree one can have only the stable type $V_3$.
The converse is immediate from stability.  The asserted vanishing
is the negative-slope section vanishing.

These arguments use the degree and slope classification on a
geometric Fargues--Fontaine curve
\cite[Proposition~II.2.10 and Theorem~II.2.14]{FS21}.
Fiberwise surjectivity is the relative subbundle condition; the
local forms in Lemma~\ref{lem:relative-section-facts} and descent
give the corresponding relative statements.
\end{proof}

Let
\[
 \widetilde X=\varprojlim_U\Spa((L^U)^\flat)
\]
be the perfected completed connected-marking tower.  This notation
does not replace the finite-stage coefficient convention of
Definition~\ref{def:stage-cochains} by cochains on a completion of
the whole tower.

\begin{proposition}[Integral frame comparison]
\label{prop:integral-frames}
There exists a determinant identification
$\iota:\det V_3\xrightarrow{\sim}\det V_2$ over $k$
for which the following conclusions hold.  Fix such an identification
for the sequel.  There is an isomorphism
\begin{equation}
 \widetilde X\simeq\operatorname{Surj}(V_3,V_2).
 \label{eq:connected-surjections}
\end{equation}
A surjection has a unique kernel frame for which its determinant
identification is $\iota$.  Under
\eqref{eq:connected-surjections}, this is an integral étale Tate
generator for the corresponding deformation.
The source-frame group $P$ acts by reduced norm on this normalization,
and a target-frame change in $J$ has its corresponding constant
determinant adjustment.
\end{proposition}

\begin{proof}
We construct the quotient frame, normalize its determinant, and then
recover the markings from an arbitrary surjection.  The reconstruction
will show that every surjection arises from an integral connected marking.

\smallskip\noindent\emph{Constructing the quotient frame.}
Over $\mathcal T_{\mathrm{et}}$, Proposition~\ref{prop:generator-ball} identifies
the chosen Tate generator with a nonzero section
$\mathcal O\to V_3$.  After adding a connected marking, there is
also a quotient frame.  To construct it, orient the formal marking
as an isomorphism from the specialized formal law to $\Gamma_2$.
Its coefficients are bounded by one on perfectoid tests.  Indeed,
its leading coefficient has norm $|t|^{1/(d-1)}<1$, and the same
coefficient comparison gives
$b_n^d-t^n b_n=Q_n(b_1,\ldots,b_{n-1})$ with integral coefficients.
If all earlier coefficients have norm at most one, a root with
$|b_n|>1$ would make the monic term strictly larger than every
other term.  Induction gives the assertion.

Consequently the formal marking can be evaluated on all small
points.  Compose it with the zeroth projection in
Lemma~\ref{lem:cover-coordinates}.  This gives an additive map
$\mathcal H_3\to\mathcal H_2$.  It is $\Qp$-linear because
multiplication by $p$ is invertible on the Honda target.  It kills
the Tate generator and hence its $\Qp$-span.  By
Lemma~\ref{lem:relative-section-facts} it comes from a unique
bundle map $V_3\to V_2$.  On a geometric fiber this map is nonzero:
near the origin the formal marking has nonzero linear coefficient,
and the zeroth projection is surjective.

Lemma~\ref{lem:stable-bundle-tests} identifies the quotient by the
Tate section with $V_2$ locally.  The induced nonzero endomorphism
of this quotient is invertible, since its endomorphism algebra is
the division algebra $D_2$.  The resulting quotient frame is
therefore an isomorphism, also relatively.  Uniqueness in relative
full faithfulness makes the construction equivariant and compatible
with all base changes.

\smallskip\noindent\emph{Normalizing the determinant.}
An initial determinant identification $\iota_0$ exists over $k$.
In the cyclic
Honda isocrystal the determinant Frobenius differs from the
rank-one degree-one isocrystal by the cyclic permutation sign
$(-1)^{i-1}$.  The possible minus sign is removed over
$\mathbb F_{p^2}$ by a Teichm\"uller unit $u$ with $u^p=-u$.
Our field contains $\mathbb F_{p^6}$, so the required identifications
are defined over $k$.  Automorphisms act on these determinant
lines by their reduced norms.

Connected markings form a $J$-torsor over $\mathcal T_{\mathrm{et}}$ and change
the quotient frames through precisely $J$.  On this torsor define
the determinant isomorphism and its discrepancy by
\begin{align*}
 \Delta(e,f)(e\wedge v_2\wedge v_3)&=f(v_2)\wedge f(v_3),\\
 r(e,f)&=\Delta(e,f)\iota_0^{-1}\in\underline{\Qp^\times}.
\end{align*}
Changing the marking by $j\in J$ multiplies $r$ by
$\Nrd(j)\in\Zp^\times$.  Consequently
$\nu=v_p(r)$ descends through this torsor to a locally constant
function $\mathcal T_{\mathrm{et}}\to\underline{\mathbb Z}$.
After extension to the field $C$ used in
Proposition~\ref{prop:generator-ball}, the space $\mathcal T_{\mathrm{et},C}$
is a punctured perfected open disc.  Nested connected closed
annuli exhaust it, and perfection preserves their topology;
thus $\nu_C$ is a single integer $m$.
The covers $S_C\to S$ constructed there also show that
$\Spa(C,C^\circ)\to\operatorname{Spd}(k)$ is $v$-surjective.
Hence $\mathcal T_{\mathrm{et},C}\to\mathcal T_{\mathrm{et}}$ is a $v$-cover and
$\nu=m$ descends globally.  Set $\iota=p^m\iota_0$.
The new discrepancy is $p^{-m}r$ and is everywhere a unit.

Each pair of connected and étale markings therefore gives an exact
sequence
\[
 0\longrightarrow\mathcal O\longrightarrow V_3
       \longrightarrow V_2\longrightarrow0
\]
whose determinant identification is a unit multiple of $\iota$.
For a fixed connected marking, multiply its étale generator by this
unit to obtain determinant equality.  There is exactly one such
generator; uniqueness glues it over $\widetilde X$.  We have thus
constructed a map from $\widetilde X$ to $\operatorname{Surj}(V_3,V_2)$,
with a determinant-normalized integral kernel frame.

\smallskip\noindent\emph{Recovering the markings.}
To prove that this map is an isomorphism, start with an arbitrary
surjection $f:V_3\twoheadrightarrow V_2$.  Determinant equality supplies
its kernel generator $e_f$.  Proposition~\ref{prop:generator-ball}
reconstructs from $e_f$ a deformation and a primitive integral
Tate generator.  Choose a connected marking locally and let $f_0$
be the quotient frame it constructs.  Write $f=b f_0$ with
$b\in D_2^\times$.  Put $r_0=\Delta(e_f,f_0)\iota^{-1}$,
which is a unit by the normalization above.  Since
$\Delta(e_f,f)=\iota$, we have $1=\Nrd(b)r_0$.  Thus
$v_p(\Nrd b)=0$, so $b\in J$ and the required change of marking
is integral.  It is unique because the marking torsor and the
integral quotient-frame torsor have the same group $J$.
Indeed, replacing the local marking by $j$ replaces $f_0$ by
$jf_0$ and $b$ by $bj^{-1}$, so the adjusted markings agree
and glue.  The resulting connected marking maps back to $f$.
Conversely, starting with a connected marking,
Proposition~\ref{prop:generator-ball} recovers its deformation from the
normalized integral generator, and the $J$-torsor uniqueness recovers
its marking.  These constructions commute with base change,
so they prove the relative isomorphism \eqref{eq:connected-surjections}.
In particular, the normalized kernel frame of every surjection is
the integral Tate generator of its reconstructed deformation.

A $p$-power rescaling of $e_f$ can change the reconstructed
deformation; the coordinate comparison still returns a primitive
basis for that deformation.  No primitivity claim about the old
deformation is needed.

\smallskip\noindent\emph{The frame actions.}
The actions can be recorded without an implicit choice of a section
of the norm map.  In the frame convention
\[
 (g,j)\cdot f=j f g^{-1},\qquad(g,j)\in P\times J,
\]
the normalized kernel frame satisfies
\begin{equation}
 e_{jfg^{-1}}=\Nrd(j)\Nrd(g)^{-1}\,g e_f.
 \label{eq:normalized-frame-action}
\end{equation}
This is the determinant identity for the exact sequence.  Reversing
all frame conventions inverts the displayed norm characters.
In either convention $H$ fixes the normalized étale frame under
transport, and the $P$ and $J$ adjustments are continuous constant
units.
\end{proof}

Put $N=\mathcal H^1(V_2^\vee)$, the relative sheaf of extensions of
$V_2$ by $\mathcal O$, with both ends framed.

\begin{lemma}[Finite marking stages of the perfected tower]
\label{lem:perfected-marking-torsor}
For each open $U\subset J$, the map
\[
 \widetilde X\longrightarrow\Spa(B^\flat),\qquad B=L^U,
\]
is a pro-étale $U$-torsor.  It is natural in $U$, in coefficient
base change, and in the commuting $P$-action.  The space
$\widetilde X$ is represented by the completed perfection of the
whole algebraic marking tower, equipped with its uniform valuation
norm; this representation is used here only to discuss its geometry.
\end{lemma}

\begin{proof}
For $V\triangleleft U$ open, $L^V/L^U$ is a finite étale
$U/V$-torsor.  Perfection commutes with finite étale base change.
After completion it remains the finite étale torsor
\[
 \Spa((L^V)^\flat)\longrightarrow\Spa(B^\flat).
\]
One may check the completed algebra explicitly in a finite field
basis: completion of its scalar extension from $B^{\mathrm{perf}}$
is $(L^V)^\flat$, since finite-dimensional valued spaces are
complete and their norms are equivalent to the coordinate norms.
This works componentwise for field products and preserves the
torsor identities.

The completed direct limit of these perfected finite-stage
algebras is uniform, perfect, and Tate, with the common
pseudouniformizer $t$.  It is therefore perfectoid.  Its associated
space is their affinoid inverse limit.  The finite-level torsor
identities pass to this limit and identify its relation with
$\underline U\times\widetilde X$.
Surjectivity can be checked on a geometric point of $\Spa(B^\flat)$:
choose compatible lifts through the surjective finite étale
covers, using compactness of their finite fibers, and extend the
resulting bounded map to the completion.  This also proves the
covering assertion.  Every construction commutes with restrictions
and the transported $P$-action.
\end{proof}

\begin{theorem}[The determinant-one two-tower comparison]
\label{thm:two-towers}
For each open $U\subset J$ and $B=L^U$ there is a natural equivalence
of $v$-stacks over $k$
\begin{equation}
 [\Spa(B^\flat)/H]\simeq[N^*/U].
 \label{eq:two-towers}
\end{equation}
It is obtained from the commuting torsor presentations
\[
 \widetilde X/U\simeq\Spa(B^\flat),\qquad
 \widetilde X/H\simeq N^*.
\]
The residual group $P/H=\Zp^\times$ acts on $N^*$ by kernel
scalars through reduced norm, up to the common inverse convention.
The $J$-action changes the quotient frame and makes the kernel
determinant adjustment in
\eqref{eq:normalized-frame-action}.  All these identifications
commute with base change and products with profinite parameter spaces.
\end{theorem}

\begin{proof}
By Proposition~\ref{prop:integral-frames}, $\widetilde X$ is the
space of surjections $V_3\to V_2$ with their uniquely normalized
kernel frames.  Quotienting the middle frame by $H$ allows the
middle bundle to vary while preserving its determinant frame.
Lemma~\ref{lem:stable-bundle-tests} identifies precisely the
resulting objects with the nowhere-split extensions of $V_2$ by
$\mathcal O$, namely $N^*$.  There is no automorphism group left
over an extension with its two ends fixed, since
$\mathcal H^0(V_2^\vee)=0$.

To see that this is a torsor presentation relatively, first choose
a local basic frame of the middle bundle.  Its determinant frame
can be corrected to the prescribed one: reduced norm
$D_3^\times\to\Qp^\times$ is surjective.  The compatible frames
form a torsor under its norm-one subgroup, which is exactly $H$,
because norm one has valuation zero.  On units, surjectivity of
$P\to\Zp^\times$ already follows from the norm of the maximal
unramified cubic subfield.  Local basic triviality then gives the
asserted $H$-torsor on the site.

Combine it with the $U$-torsor in
Lemma~\ref{lem:perfected-marking-torsor}.  The actions commute, so
the two iterated quotient presentations give
\eqref{eq:two-towers}.  Tracking determinants gives the claimed
actions.  For example, in the frame convention of
\eqref{eq:normalized-frame-action}, the residual $g\in P$ acts
on extension classes by the scalar $\Nrd(g)$; a target change
$j$ acts by pullback along $j^{-1}$ together with kernel scalar
$\Nrd(j)^{-1}$.  This states the actual action and does not choose
a group-theoretic splitting of reduced norm.  Naturality and the
parameter assertion follow from the natural torsor maps and the
uniform coordinate comparison.
\end{proof}

\subsection{Descent of the integral marking}

The normalized generator has so far been constructed on a completed
perfected tower.  We next descend each of its finite étale levels to
the algebraic marking algebra $L$.  The following elementary lemma
is the completion statement needed for this purpose.

\begin{lemma}[Finite étale sections descend]
\label{lem:finite-etale-descent}
Let $L=\colim_i B_i$ be an injective union of finite étale
$K$-algebras, with their maximum valuation norms, and let
$C=\widehat{L^{\mathrm{perf}}}$.  For every finite étale
$L$-algebra $E$, the map
\[
 \Hom_{L\text{-alg}}(E,L)\longrightarrow
 \Hom_{L\text{-alg}}(E,C)
\]
is bijective.  Each section on the right is defined at a finite
separable stage.  The assertion allows unbounded numbers of field
factors among the $B_i$.
\end{lemma}

\begin{proof}
Each $B_i^{1/p^r}$ is a finite product of complete valued fields.
First let $e\in C$ be an idempotent.  Approximate it by
$b\in B_i^{1/p^r}$ with $\lVert b-e\rVert<1$.
Then $\lVert b^2-b\rVert<1$ and $2b-1$ is a unit with inverse
of norm at most one.  Hensel's lemma in that complete finite
product gives an idempotent $e_i$ at distance less than one from
$e$.  Two such idempotents are equal: both $e(1-e_i)$ and
$e_i(1-e)$ are idempotents of norm less than one, hence zero.
Purely inseparable extensions create no idempotents, so $e=e_i$
belongs to $B_i$.  Finitely many idempotents descend to a common
stage.

The algebra $E$ descends to a finite étale algebra over some $B_i$.
After a finite idempotent partition it has a monogenic étale
presentation; each component field of $B_i$ is infinite, so the
primitive-element argument applies also to its finite étale
products.  Descend the idempotents selected by the proposed
section.  It remains to descend finitely many simple roots.

Let $f(z)=0$ with $f\in B_i[T]$ and $f'(z)$ invertible in $C$.
Approximate $z$ and $f'(z)^{-1}$ simultaneously by $b,c$ in a
single $B_j^{1/p^r}$, so closely that
\[
 \lVert1-cf'(b)\rVert<1.
\]
The geometric series in this complete algebra makes $f'(b)$
invertible with inverse norm at most $\lVert c\rVert$.
Let $M\geq1$ bound the quadratic Taylor remainder of $f$ on
the unit ball about $b$, and put $C_0=\max(1,\lVert c\rVert)$.
The approximations may also be chosen so that
\[
 \lVert f(b)\rVert C_0^2M<1.
\]
Newton iteration remains in $B_j^{1/p^r}$, has uniformly bounded
inverse derivatives, and converges there to a root $z_j$.
Choose the initial approximation inside a simple-root uniqueness
ball about $z$; the divided difference is a unit on that ball,
so $z_j=z$.  Componentwise, this root is separable over $B_j$
but belongs to its purely inseparable extension, and hence was
already in $B_j$.  All the finitely many roots use a common
stage.  This proves surjectivity, and injectivity follows from
$L\hookrightarrow C$.
\end{proof}

\begin{lemma}[The normalized algebraic étale marking]
\label{lem:etale-marking}
The étale quotient of $G_L$ has a compatible integral Tate
generator defined over $L$, obtained by determinant normalization.
It is fixed by $H$.  Transport by $P$ changes it by the constant
reduced-norm unit, up to the common inverse convention; transport
by $J$ has the constant determinant adjustment determined by
\eqref{eq:normalized-frame-action}.  Each finite torsion level
of this marking is defined at a finite separable marking stage.
\end{lemma}

\begin{proof}
The integral Tate frame constructed in
Proposition~\ref{prop:integral-frames} gives, at each level $l$,
a section of the finite étale generator algebra of
$G_L^{\mathrm{et}}[p^l]$ after passage to
$\widehat{L^{\mathrm{perf}}}$.
Lemma~\ref{lem:finite-etale-descent} descends this section to
some finite separable marking stage.  The descended sections
are compatible: their transition identities hold in the
completion and hence in $L$ by injectivity.  The same argument
descends the action identities, including the $H$-invariance.
The finite stage can depend on $l$; no single finite stage
containing the entire Tate generator is required.
\end{proof}

\begin{lemma}[The generator-lift torsors]
\label{lem:lift-torsors}
The scheme of lifts in $G_L[p^l]$ of the normalized étale
generator has coordinate algebra
\begin{equation}
 R_l=L^{1/p^{2l}}=L^{1/d^l}\subset R=L^{\mathrm{perf}}.
 \label{eq:lift-rings}
\end{equation}
It is a torsor under the marked group $\Gamma_2[p^l]_L$.
The transition maps give the indicated inclusions inside $R$.
The action of $H$ commutes with translation by this constant
group.  The other frame actions conjugate translations by
continuous constant Honda automorphisms and scalar norm units.
All coactions and their finite-level action identities are
defined and continuous after passage to a sufficiently large
finite separable stage.
\end{lemma}

\begin{proof}
The fiber over the marked étale generator is a torsor for the
connected kernel.  The connected marking identifies that
kernel with $\Gamma_2[p^l]$.
By Lemma~\ref{lem:primitive-coordinates}, its algebra before
base change along the chosen generator is
\[
 k[[s_l]][1/t]
 \quad\text{over}\quad
 k[[s_l^{d^l}]][1/t],
\]
and this is the full lift algebra of rank $d^l$.
The generator marking of Lemma~\ref{lem:etale-marking}
gives its base change to $L$.

Every finite separable extension of $K$ has one-element
$p$-basis $t$, and the same is true of their ind-étale union.
A finite purely inseparable extension of degree $d^l$ in
its perfection is therefore $L^{1/d^l}$.  The displayed lift
algebra retains that degree after the separable base change:
componentwise, $s_l^{d^l}$ is a $p$-basis of its Laurent-series
field and remains a $p$-basis after any separable extension.
This proves \eqref{eq:lift-rings}, also for products by a
common finite-stage argument.  Purely inseparable embeddings
are unique, so compatibility of the marked generators
identifies the transition maps with the inclusions inside $R$.

The group $H$ fixes both the connected frame and the normalized
étale generator.  Its transported action on each lift torsor
therefore commutes with translations in the named Honda group.
For the remaining actions, transport followed by restoration
of the normalized étale generator rescales by its constant
norm unit; changing the connected frame also applies the named
Honda automorphism.  These are continuous constant group
automorphisms and satisfy the action identities because the
markings do.  At a fixed level all these statements involve
finite torsion algebras and finitely many coefficients.  Those
coefficients lie in a finite separable stage, and the resulting
maps between complete finite-dimensional valued spaces are
continuous.  Enlarging that stage handles any specified finite
list of levels and identities.

For later use, the order-two action has a direct sign check.
The central elements $\tau=-1\in P$ and $j=-1\in J$ both act
on surjections by $f\mapsto-f$.  Since $V_2$ has rank two,
$\Delta(e,-f)=\Delta(e,f)$, so the normalized generator $e$,
its reconstructed deformation, and its étale frame are unchanged.
The connected marking $\phi$ is replaced by $-\phi$.
A translation labeled by $c\in\Gamma_2[p^l]$ uses
$\phi^{-1}(c)$; after this change it uses $\phi^{-1}(-c)$.
Thus $\tau$, whose reduced norm is $-1$, conjugates the named
translation group by $[-1]$.  Its action on the one-dimensional
tangent of the Cartier dual is $-1$, as required in
Lemma~\ref{lem:distribution-operator}.
\end{proof}

\subsection{The negative period space}

The final geometric description will be used to compute cohomology.
It requires a geometric untilt only after the space and its frame
actions have been defined over $k$.

\begin{lemma}[An untilt presentation of the negative space]
\label{lem:negative-presentation}
Let $C^\flat$ be an algebraically closed perfectoid field of
characteristic $p$ containing $k$, choose an algebraically closed
untilt $C/\mathbb C_p$, and let $F/\Qp$ be unramified quadratic.
After choosing its embedding at the untilt divisor there are
isomorphisms
\begin{equation}
 N_{C^\flat}\simeq
       (\mathbb G_{a,C})^\diamond/\underline F,
 \qquad
 (N^*)_{C^\flat}\simeq
       [ (\mathbb A^1_C\setminus\underline F)^\diamond/
          \underline F].
 \label{eq:negative-presentation}
\end{equation}
Here $\mathbb A^1_C$ is the analytic affine line.  In the second
expression the $F$-support is removed fiberwise
before taking the translation quotient.  The action of
$F^\times$ through the quotient frame is multiplication in this
presentation, up to inversion and the choice of unramified
embedding.  The additional kernel determinant adjustment
from Theorem~\ref{thm:two-towers} is a scalar in $\Zp^\times$.
\end{lemma}

\begin{proof}
On the Fargues--Fontaine curve with coefficient field $F$,
the chosen untilt divisor gives \citep[Proposition~II.2.3]{FS21}
\[
 0\longrightarrow\mathcal O_F(-1)
  \longrightarrow\mathcal O_F
  \longrightarrow i_*C\longrightarrow0.
\]
The degree-zero section sheaf of $\mathcal O_F$ is
$\underline F$, its higher relative cohomology sheaves vanish,
and the negative line has no sections.  The cohomology sequence
therefore identifies
\[
 \mathcal H^1(\mathcal O_F(-1))
       \simeq (\mathbb G_{a,C})^\diamond/\underline F.
\]
Finite pushforward along the unramified coefficient extension
takes $\mathcal O_F(-1)$ to $V_2^\vee$, giving the first
isomorphism.  The preimage of the zero extension class is
exactly the sheaf $\underline F$.  Its complement consists
precisely of classes nonzero on every geometric fiber, proving
the second isomorphism as an open subquotient.

The sequence is equivariant for multiplication by $F^\times$.
The embedding and frame conventions may conjugate or invert
this action.  The independent determinant normalization of
the kernel contributes exactly the scalar described in
Theorem~\ref{thm:two-towers}.  For later use, a scalar
$u\in\Zp^\times$ has norm $u^2$ in $F/\Qp$, which is
trivial modulo $3$; thus this adjustment does not change the
nontrivial residue-norm character on $\mathcal O_F^\times$.
\end{proof}

The two-tower comparison and its normalized scalar action now supply the
geometric input for the completed calculation. Before using them, we
establish the finite analytic resolution that controls continuous cochains.
